\documentclass[11pt,a4paper]{article}

% ─────────────────────────────────────────────────────────────────────────────
%  Metodología para la Reconstrucción de Matrices Insumo-Producto
%  Argentina · Brasil · México · Uruguay
%  Proyecto matrices-mip (edwinpulgarin/matrices-mip)
%  Notación alineada al marco CEPAL "MIP Extendida Ambientalmente" (2025).
% ─────────────────────────────────────────────────────────────────────────────

\usepackage[utf8]{inputenc}
\usepackage[T1]{fontenc}
\usepackage[spanish,es-tabla]{babel}
\usepackage{amsmath,amssymb,amsfonts}
\usepackage{booktabs}
\usepackage{longtable}
\usepackage{array}
\usepackage{geometry}
\geometry{margin=2.5cm}
\usepackage{enumitem}
\usepackage{xcolor}
\usepackage{graphicx}
\usepackage{caption}
\usepackage[hidelinks]{hyperref}
\usepackage{fancyhdr}

\hypersetup{
  colorlinks=true,
  linkcolor=blue!45!black,
  urlcolor=blue!50!black,
  citecolor=blue!45!black,
}

\pagestyle{fancy}
\fancyhf{}
\rhead{\small Metodología MIP — Argentina, Brasil, México, Uruguay}
\lfoot{\small Proyecto \texttt{matrices-mip}}
\rfoot{\small \thepage}
\renewcommand{\headrulewidth}{0.4pt}

\newcommand{\R}{\mathbb{R}}
\newcommand{\diag}{\operatorname{diag}}
\newcommand{\matname}[1]{\textbf{#1}}

\title{\textbf{Metodología para la Reconstrucción de Matrices\\
Insumo-Producto a partir de Cuadros de Oferta y Utilización}\\[4pt]
\large Argentina, Brasil, México y Uruguay}
\author{Proyecto \texttt{matrices-mip}}
\date{Corte de datos: junio de 2026}

\begin{document}

\maketitle
\thispagestyle{fancy}

\begin{abstract}
\noindent
Este documento describe la metodología empleada para construir y reconstruir
una serie de \textbf{36 matrices insumo-producto (MIP)} anuales para cuatro
economías latinoamericanas —Argentina, Brasil, México y Uruguay—, así como las
fuentes, supuestos, validaciones y límites del trabajo. La serie combina
matrices \emph{directas} (publicadas por las oficinas nacionales de estadística)
con matrices \emph{reconstruidas} a partir de Cuadros de Oferta y Utilización
(COU/SUT) bajo el supuesto de tecnología de industria. La notación sigue el
marco de la CEPAL para la \emph{Matriz Insumo-Producto Extendida Ambientalmente
y Huella de Carbono} (2025), de modo que la serie quede lista para una eventual
extensión ambiental. Se documentan el modelo de Leontief (encadenamientos hacia
atrás), el modelo de Ghosh (encadenamientos hacia adelante), los multiplicadores
de producción, la validación matemática y la auditoría de cobertura sectorial.
\end{abstract}

\tableofcontents
\newpage

% ═════════════════════════════════════════════════════════════════════════════
\section{Objetivos y Alcance}
% ═════════════════════════════════════════════════════════════════════════════

\subsection{Objetivo general}
Construir, documentar y validar una serie comparable de matrices insumo-producto
(MIP) anuales para Argentina, Brasil, México y Uruguay, preservando la
trazabilidad completa entre cada matriz publicada, su fuente original y las
transformaciones aplicadas.

\subsection{Objetivos específicos}
\begin{enumerate}[leftmargin=1.6em]
  \item Reconstruir matrices industria por industria a partir de los Cuadros de
        Oferta y Utilización (COU) cuando la oficina nacional no publica una MIP
        simétrica para el año requerido.
  \item Incorporar las MIP directas publicadas (México, Uruguay~2016,
        Argentina~1997) sin alterar sus cifras oficiales, adjuntando el COU de
        referencia cuando exista.
  \item Calcular los coeficientes técnicos, la inversa de Leontief, los
        coeficientes de Ghosh y los multiplicadores de producción de cada matriz.
  \item Validar la consistencia matemática de cada MIP y auditar que ninguna
        actividad económica de la fuente quede sin representación en la matriz
        final.
  \item Alinear la notación al marco metodológico de la CEPAL para dejar la serie
        preparada para una extensión ambiental (huella de carbono) futura.
\end{enumerate}

\subsection{Alcance}

\paragraph{Alcance temporal y geográfico.}
La serie cubre cuatro países y abarca desde 1997 (Argentina) hasta 2022
(Argentina), con la cobertura por país descrita en el Cuadro~\ref{tab:inventario}.

\paragraph{Alcance sectorial.}
El nivel de desagregación lo fija cada fuente nacional y varía por año: Brasil a
51 actividades (base 2000) o 68 (nivel SCN/IBGE 2010--2021); Argentina entre 107
y 162; México entre 20 (2003) y 262--263 (rama SCIAN); Uruguay entre 95 y 128. No
se fuerza una clasificación común entre países: cada matriz conserva la
desagregación de su fuente.

\paragraph{Alcance de las fuentes.}
Todas las fuentes son oficiales: INDEC (Argentina), IBGE (Brasil), INEGI
(México) y BCU (Uruguay), complementadas por el repositorio COU/MIP de la CEPAL.
El detalle por país está en la Sección~\ref{sec:fuentes}.

\paragraph{Lo que este trabajo NO incluye (aún).}
La extensión ambiental (matrices de presiones, intensidades de emisión y huella
de carbono) se describe como marco de referencia en la
Sección~\ref{sec:ambiental} pero \emph{no} se estima, porque todavía no se
dispone de vectores de emisiones por sector para estos cuatro países.

% ═════════════════════════════════════════════════════════════════════════════
\section{Conceptos Básicos y Notación}
% ═════════════════════════════════════════════════════════════════════════════

Sea $n$ el número de actividades económicas de una economía. La MIP describe los
flujos monetarios entre actividades para un año dado.

\subsection{Estructura de la matriz insumo-producto}

\paragraph{Matriz de consumo intermedio.}
\begin{equation}
Z =
\begin{bmatrix}
z_{11} & z_{12} & \cdots & z_{1n}\\
z_{21} & z_{22} & \cdots & z_{2n}\\
\vdots & \vdots & \ddots & \vdots\\
z_{n1} & z_{n2} & \cdots & z_{nn}
\end{bmatrix}
\in \R^{n\times n},
\end{equation}
donde $z_{ij}$ es el valor del producto de la actividad $i$ usado como insumo por
la actividad $j$.

\paragraph{Vectores económicos.}
El vector de \textbf{producción bruta} $x\in\R^{n}$, la \textbf{demanda final}
$y\in\R^{n}$ y el \textbf{valor agregado} $v\in\R^{n}$ satisfacen
\begin{equation}
y_i = C_i + G_i + I_i + X_i - M_i,
\end{equation}
con $C$ consumo de hogares, $G$ consumo de gobierno, $I$ formación bruta de
capital, $X$ exportaciones y $M$ importaciones. La matriz diagonal de producción
es $\hat{X}=\diag(x)$.

\subsection{Balance básico}
El equilibrio contable de la economía establece que la producción de cada
actividad iguala su uso intermedio más su demanda final:
\begin{equation}
x_i = \sum_{j=1}^{n} z_{ij} + y_i, \qquad\text{o en forma matricial}\qquad
x = Z\,\mathbf{1} + y,
\label{eq:balance}
\end{equation}
donde $\mathbf{1}=[1,\dots,1]^{T}$. Por columnas, la producción iguala el
consumo intermedio más el valor agregado: $x_j=\sum_i z_{ij}+v_j$.

\subsection{Modelo de Leontief: encadenamientos hacia atrás}

\paragraph{Coeficientes técnicos.}
\begin{equation}
a_{ij}=\frac{z_{ij}}{x_j}, \qquad A = Z\,\hat{X}^{-1}\in\R^{n\times n}.
\end{equation}
El coeficiente $a_{ij}$ indica cuántas unidades monetarias del producto $i$ se
requieren para producir una unidad monetaria del producto $j$.

\paragraph{Ecuación fundamental e inversa de Leontief.}
Sustituyendo $Z=A\hat{X}$ en \eqref{eq:balance} se obtiene $x = Ax + y$, de donde
\begin{equation}
(I-A)\,x = y \qquad\Longrightarrow\qquad x = (I-A)^{-1}y = L\,y,
\end{equation}
\begin{equation}
L=(I-A)^{-1}\in\R^{n\times n}.
\end{equation}
El elemento $l_{ij}$ representa la producción total —directa e indirecta— del
sector $i$ necesaria para satisfacer una unidad de demanda final del sector $j$.

\subsection{Modelo de Ghosh: encadenamientos hacia adelante}

\paragraph{Coeficientes de distribución.}
\begin{equation}
b_{ij}=\frac{z_{ij}}{x_i}, \qquad B = \hat{X}^{-1}Z\in\R^{n\times n},
\end{equation}
donde $b_{ij}$ es la proporción de la producción del sector $i$ que se destina al
sector $j$ como insumo.

\paragraph{Inversa de Ghosh.}
\begin{equation}
G=(I-B)^{-1}\in\R^{n\times n}.
\end{equation}
$g_{ij}$ muestra cómo un cambio en la oferta primaria del sector $i$ se propaga
hacia adelante por toda la cadena productiva.

\subsection{Multiplicadores y encadenamientos}

\paragraph{Multiplicador de producción.}
\begin{equation}
MP_j = \sum_{i=1}^{n} l_{ij},
\end{equation}
mide el efecto total sobre la producción de la economía ante un aumento unitario
de la demanda final del sector $j$.

\paragraph{Índices de encadenamiento.}
\begin{equation}
BL_j = \frac{\sum_{i} l_{ij}}{\tfrac{1}{n}\sum_{i,j} l_{ij}},
\qquad
FL_i = \frac{\sum_{j} g_{ij}}{\tfrac{1}{n}\sum_{i,j} g_{ij}}.
\end{equation}
$BL_j$ (backward linkage) mide el poder de arrastre del sector como demandante;
$FL_i$ (forward linkage) mide su poder de empuje como oferente. La clasificación
de sectores resulta de cruzar ambos índices:
\begin{itemize}[leftmargin=1.6em,nosep]
  \item \textbf{Clave}: $BL>1$ y $FL>1$.
  \item \textbf{Estratégico hacia atrás}: $BL>1$ y $FL<1$.
  \item \textbf{Estratégico hacia adelante}: $BL<1$ y $FL>1$.
  \item \textbf{Independiente}: $BL<1$ y $FL<1$.
\end{itemize}

\subsection{Resumen de notación}
\begin{table}[h]
\centering
\small
\begin{tabular}{cl}
\toprule
\textbf{Símbolo} & \textbf{Descripción}\\
\midrule
$n$        & Número de actividades económicas\\
$Z$        & Matriz de consumo intermedio ($n\times n$)\\
$x$        & Vector de producción bruta ($n\times 1$)\\
$y$        & Vector de demanda final ($n\times 1$)\\
$v$        & Vector de valor agregado ($n\times 1$)\\
$\hat{X}$  & Matriz diagonal de producción, $\diag(x)$\\
$A$        & Coeficientes técnicos, $A=Z\hat{X}^{-1}$\\
$L$        & Inversa de Leontief, $L=(I-A)^{-1}$\\
$B$        & Coeficientes de distribución, $B=\hat{X}^{-1}Z$\\
$G$        & Inversa de Ghosh, $G=(I-B)^{-1}$\\
$MP_j$     & Multiplicador de producción del sector $j$\\
$BL,\,FL$  & Encadenamientos hacia atrás / adelante\\
\midrule
\multicolumn{2}{l}{\emph{Objetos del COU (reconstrucción)}}\\
$V$        & Tabla de oferta/producción (industria $\times$ producto)\\
$U$        & Tabla de utilización intermedia (producto $\times$ industria)\\
$q$        & Producción por producto\\
$g$        & Producción bruta por industria\\
$D$        & Matriz de participación de mercado (\emph{market share})\\
\bottomrule
\end{tabular}
\caption{Notación utilizada en el documento.}
\end{table}

% ═════════════════════════════════════════════════════════════════════════════
\section{El Cuadro de Oferta y Utilización (COU)}
% ═════════════════════════════════════════════════════════════════════════════

La mayoría de las matrices de esta serie se reconstruye a partir de Cuadros de
Oferta y Utilización (COU; en inglés \emph{Supply-Use Tables}, SUT). El COU es el
insumo primario del Sistema de Cuentas Nacionales y, a diferencia de la MIP,
distingue \textbf{productos} (filas) de \textbf{industrias/actividades}
(columnas).

\subsection{Tablas de oferta y utilización}
\begin{itemize}[leftmargin=1.6em]
  \item \textbf{Tabla de oferta} $V$ (industria $\times$ producto): registra qué
        produce cada industria. La suma por producto da la producción
        $q=\mathbf{1}^{T}V$.
  \item \textbf{Tabla de utilización} $U$ (producto $\times$ industria): registra
        qué productos consume cada industria como insumo intermedio, más la
        demanda final por componentes.
\end{itemize}

\subsection{Productos versus industrias}
Un punto metodológico delicado es no confundir productos con sectores. Un
producto puede aparecer en la cobertura del COU y no como sector de la MIP final,
porque la transformación COU$\,\to\,$MIP pasa por la matriz de participación de
mercado $D$ (Sección~\ref{sec:reconstruccion}).

\subsection{Valoración y origen}
El COU plantea dos retos de valoración que el procedimiento debe resolver:
\begin{enumerate}[leftmargin=1.6em]
  \item \textbf{Precios básicos vs.\ precios de comprador.} La utilización suele
        publicarse a precios de comprador (incluye márgenes de comercio y
        transporte e impuestos netos sobre los productos), mientras que la
        producción está a precios básicos. La MIP doméstica se construye a
        precios básicos.
  \item \textbf{Origen doméstico vs.\ importado.} La MIP doméstica solo contempla
        la tecnología y producción nacional. El uso intermedio importado se
        separa para que la demanda final quede satisfecha exclusivamente por
        producción local.
\end{enumerate}

% ═════════════════════════════════════════════════════════════════════════════
\section{Metodología de Reconstrucción COU \texorpdfstring{$\to$}{->} MIP}
\label{sec:reconstruccion}
% ═════════════════════════════════════════════════════════════════════════════

\subsection{Supuesto de tecnología de industria}
La transformación de un sistema rectangular oferta-utilización a una MIP
cuadrada industria por industria emplea el \textbf{supuesto de tecnología de
industria} (\emph{industry-technology assumption}), que asume que cada industria
usa la misma estructura de insumos para todos los productos que fabrica. Este
supuesto preserva la no negatividad de los coeficientes.

\subsection{Matriz de participación de mercado y derivación de \texorpdfstring{$Z$}{Z}}
Se define la matriz de participación de mercado $D$ a partir de la oferta:
\begin{equation}
D = V\,\hat{q}^{-1}, \qquad \hat{q}=\diag(q),
\end{equation}
donde $d_{kp}$ es la proporción del producto $p$ fabricada por la industria $k$.
El consumo intermedio industria por industria se obtiene proyectando la
utilización nacional por la participación de mercado:
\begin{equation}
Z = D\,U_{\text{nacional}}, \qquad
A = Z\,\hat{g}^{-1},\quad L=(I-A)^{-1},
\end{equation}
\begin{equation}
B=\hat{g}^{-1}Z,\quad G=(I-B)^{-1},\qquad
f = D\,y_{\text{doméstica}},
\end{equation}
con $g$ la producción bruta por industria y $f$ la demanda final ya expresada por
industria.

\subsection{Tratamiento de precios y de las importaciones}
Cuando la fuente publica la utilización a precios de comprador, $U$ se reescala
por producto a precios básicos con un factor
$\phi_p = q^{\text{pb}}_p / (\text{oferta total}_p - \text{importaciones}_p)$,
acotado a $\phi_p\le 1$ por la identidad \emph{precio básico} $\le$ \emph{precio
de comprador}. El uso importado se aísla en $U_{\text{importada}}$ y no entra en
la MIP doméstica.

\subsection{Cierre de la demanda final y conciliación menor}
Para que se cumpla oferta $=$ demanda a precios básicos, la demanda final
doméstica se cierra como residual $y = q^{\text{pb}} - U\mathbf{1}$ cuando la
fuente no permite una conciliación exacta. La política de conciliación es
conservadora:
\begin{itemize}[leftmargin=1.6em,nosep]
  \item Primero se revisan parsers, totales, códigos y puentes de precios.
  \item No se alteran cifras para «forzar» resultados.
  \item Solo se aplica una conciliación menor tipo RAS cuando el desbalance es
        pequeño, localizado y justificable como cierre o redondeo (por ejemplo,
        Brasil 2000--2009 y Uruguay~2016).
  \item Toda conciliación deja traza en las hojas \texttt{ajuste\_cierre} y
        \texttt{Z\_pre\_conciliacion} del libro Excel.
\end{itemize}

\subsection{Sectores con diagonal cero}
Una actividad con $z_{ii}=0$ \emph{no} se elimina automáticamente. Si el sector
registra producción, valor agregado, ventas, compras o demanda final, se
conserva; la diagonal nula solo indica ausencia de autoconsumo. Estos casos se
documentan como \texttt{diagonal\_cero\_con\_flujos} (por ejemplo, Servicios
Domésticos en Argentina).

\subsection{Resumen del procedimiento}
\begin{enumerate}[leftmargin=1.6em,nosep]
  \item Leer y normalizar el COU de la fuente ($V$, $U$, demanda final, VA).
  \item Separar uso nacional de importado y llevar $U$ a precios básicos.
  \item Calcular $D=V\hat{q}^{-1}$ y $Z=D\,U_{\text{nacional}}$.
  \item Derivar $A,L,B,G$ y cerrar la demanda final.
  \item Verificar la inversa de Leontief y los multiplicadores.
  \item Empaquetar el resultado en Excel con todas las hojas de trazabilidad.
\end{enumerate}

% ═════════════════════════════════════════════════════════════════════════════
\section{Fuentes de Datos por País}
\label{sec:fuentes}
% ═════════════════════════════════════════════════════════════════════════════

\paragraph{Argentina (INDEC / CEPAL).}
MIP directa de 1997 (MIPAr97). COU de 2004 (año base) y 2018--2022. Los años
2005--2017 \emph{no} están disponibles como COU publicado: el repositorio CEPAL
solo ofrece 2004, 2018, 2019, 2020 y 2021, e INDEC añade 2022. Ese hueco no puede
llenarse con fuente pública.

\paragraph{Brasil (IBGE / CEPAL).}
Serie reconstruida 2000--2021. Para 2000--2009 se usa el COU de la CEPAL con base
2000; para 2010--2021, las Tablas de Recursos e Usos (TRU) del IBGE a nivel 68.

\paragraph{México (INEGI / CEPAL).}
MIP directas de 2003, 2008, 2013 y 2018, publicadas por INEGI/CEPAL al nivel de
rama SCIAN. A 2008 y 2013 se les adjunta su COU matricial de referencia; 2003 y
2018 no disponen de un COU compatible por rama.

\paragraph{Uruguay (BCU / CEPAL).}
MIP directa de 2016 (BCU). Reconstrucciones desde COU detallado para 2012 y 2017.
El COU 2012, publicado por el BCU en una hoja única que apila oferta, utilización
y valor agregado (134 productos $\times$ 107 industrias), se reconstruye con el
mismo procedimiento que 2017.

\medskip
\noindent\textbf{Trazabilidad de fuentes.} Al agregar una fuente nueva se
actualizan simultáneamente \texttt{FUENTES\_EXTERNAS\_HISTORICO.md}, el handoff
del proyecto, y las hojas \texttt{fuente\_resumen} y \texttt{fuente\_notas} de los
libros Excel regenerados.

% ═════════════════════════════════════════════════════════════════════════════
\section{Matrices Directas y COU de Referencia}
% ═════════════════════════════════════════════════════════════════════════════

Las MIP directas se publican tal cual, sin recálculo de sus celdas. Cuando existe
el COU del mismo marco, se adjunta como \emph{referencia} dentro del mismo libro
Excel (hojas \texttt{src\_*}), marcando el tipo
\texttt{MIP\_directa\_con\_COU\_referencia}, sin relabelar la matriz como
reconstruida y sin tocar la auditoría. Así, cada MIP lleva su COU o, en su
defecto, una nota de fuente detallada que explica por qué no lo lleva.

\medskip
\noindent Cobertura de COU sobre las 36 matrices:
\begin{itemize}[leftmargin=1.6em,nosep]
  \item \textbf{30 reconstruidas} desde COU (el COU es la fuente y se adjunta).
  \item \textbf{3 directas con COU de referencia}: México~2008, México~2013,
        Uruguay~2016.
  \item \textbf{3 directas sin COU compatible}, con fuente explicada al inicio:
        México~2003, México~2018, Argentina~1997.
\end{itemize}

% ═════════════════════════════════════════════════════════════════════════════
\section{Validación y Auditoría}
% ═════════════════════════════════════════════════════════════════════════════

\subsection{Validación matemática}
Cada MIP pasa un conjunto de pruebas estructurales: matrices cuadradas y
etiquetas alineadas ($Z,A,L$), no negatividad de $Z,A,g$, consistencia de
$A-Z\hat{g}^{-1}$, reproducción de la inversa de Leontief y de Ghosh, conteo de
celdas negativas y de sectores con demanda final o valor agregado residual
negativo.

\subsection{Auditoría de cobertura sectorial}
Una auditoría independiente verifica que toda actividad presente en la fuente
quede representada como sector en la MIP final (o quede documentada su exclusión
con justificación). El objetivo operativo es \emph{cero} actividades de fuente sin
incorporar.

\subsection{Alertas metodológicas}
Algunas pruebas son diagnósticas: pueden señalar avisos económicos (por ejemplo,
demanda final negativa en ciertos productos) sin invalidar el álgebra matricial.
El caso más sensible es Uruguay, donde el cierre residual a precios básicos genera
demandas finales negativas en algunos productos; se mantiene la alerta en lugar de
conciliar cifras materiales.

% ═════════════════════════════════════════════════════════════════════════════
\section{Inventario de Matrices}
% ═════════════════════════════════════════════════════════════════════════════

El Cuadro~\ref{tab:inventario} lista las 36 matrices de la serie con su tipo,
nivel sectorial y fuente.

\begin{center}
\small
\begin{longtable}{llrll}
\toprule
\textbf{País} & \textbf{Año} & \textbf{Sect.} & \textbf{Tipo} & \textbf{Fuente} \\
\midrule
\endhead
Argentina & 1997 & 124 & Directa & MIPAr97 INDEC \\
Argentina & 2004 & 162 & Reconstruida & COU INDEC/CEPAL \\
Argentina & 2018 & 107 & Reconstruida & COU INDEC/CEPAL \\
Argentina & 2019 & 107 & Reconstruida & COU INDEC/CEPAL \\
Argentina & 2020 & 107 & Reconstruida & COU INDEC/CEPAL \\
Argentina & 2021 & 107 & Reconstruida & COU INDEC/CEPAL \\
Argentina & 2022 & 107 & Reconstruida & COU INDEC/CEPAL \\
Brasil & 2000 & 51 & Reconstruida & COU CEPAL base 2000 \\
Brasil & 2001 & 51 & Reconstruida & COU CEPAL base 2000 \\
Brasil & 2002 & 51 & Reconstruida & COU CEPAL base 2000 \\
Brasil & 2003 & 51 & Reconstruida & COU CEPAL base 2000 \\
Brasil & 2004 & 51 & Reconstruida & COU CEPAL base 2000 \\
Brasil & 2005 & 51 & Reconstruida & COU CEPAL base 2000 \\
Brasil & 2006 & 51 & Reconstruida & COU CEPAL base 2000 \\
Brasil & 2007 & 51 & Reconstruida & COU CEPAL base 2000 \\
Brasil & 2008 & 51 & Reconstruida & COU CEPAL base 2000 \\
Brasil & 2009 & 51 & Reconstruida & COU CEPAL base 2000 \\
Brasil & 2010 & 68 & Reconstruida & TRU IBGE niv.68 \\
Brasil & 2011 & 68 & Reconstruida & TRU IBGE niv.68 \\
Brasil & 2012 & 68 & Reconstruida & TRU IBGE niv.68 \\
Brasil & 2013 & 68 & Reconstruida & TRU IBGE niv.68 \\
Brasil & 2014 & 68 & Reconstruida & TRU IBGE niv.68 \\
Brasil & 2015 & 68 & Reconstruida & TRU IBGE niv.68 \\
Brasil & 2016 & 68 & Reconstruida & TRU IBGE niv.68 \\
Brasil & 2017 & 68 & Reconstruida & TRU IBGE niv.68 \\
Brasil & 2018 & 68 & Reconstruida & TRU IBGE niv.68 \\
Brasil & 2019 & 68 & Reconstruida & TRU IBGE niv.68 \\
Brasil & 2020 & 68 & Reconstruida & TRU IBGE niv.68 \\
Brasil & 2021 & 68 & Reconstruida & TRU IBGE niv.68 \\
México & 2003 & 20 & Directa & MIP INEGI/CEPAL \\
México & 2008 & 262 & Directa (+COU ref.) & MIP INEGI/CEPAL \\
México & 2013 & 262 & Directa (+COU ref.) & MIP INEGI/CEPAL \\
México & 2018 & 263 & Directa & MIP INEGI/CEPAL \\
Uruguay & 2012 & 107 & Reconstruida & COU detallado BCU \\
Uruguay & 2016 & 128 & Directa (+COU ref.) & MIP directa BCU \\
Uruguay & 2017 & 95 & Reconstruida & COU CEPAL/BCU \\
\bottomrule
\caption{Inventario de las 36 matrices insumo-producto de la serie. Validación
estructural: 36/36 OK; auditoría: 0 sectores de fuente pendientes por
incorporar.}
\label{tab:inventario}
\end{longtable}
\end{center}

% ═════════════════════════════════════════════════════════════════════════════
\section{Marco de Extensión Ambiental (referencia)}
\label{sec:ambiental}
% ═════════════════════════════════════════════════════════════════════════════

La serie está preparada para una extensión ambiental siguiendo el marco de la
CEPAL. Esta sección documenta ese marco como referencia; \textbf{no} corresponde a
resultados ya calculados.

\paragraph{Presiones e intensidades.}
Dada una matriz de presiones ambientales $D_1\in\R^{k\times n}$ (con $k$
indicadores, por ejemplo emisiones de CO\textsubscript{2}, CH\textsubscript{4} y
N\textsubscript{2}O en tCO\textsubscript{2}eq), los coeficientes de intensidad
directa son $D = D_1\hat{X}^{-1}$ y los multiplicadores ambientales totales
$D_a = D\,L$.

\paragraph{Huella de carbono.}
Con $\alpha_1=\diag(\gamma_1/x)$ la intensidad de emisiones por sector y
$m=\alpha_1 L$ el vector de multiplicadores de emisiones, la huella de carbono de
la demanda final es
\begin{equation}
CF = m\cdot y = \sum_{j} m_j\, y_j,
\end{equation}
descomponible por componente de demanda ($CF_C, CF_G, CF_I, CF_X$) y por sector.

\paragraph{Requisito pendiente.}
Implementar esta extensión exige vectores de emisiones por sector
($\gamma$) para Argentina, Brasil, México y Uruguay, que aún no se han
incorporado. Para Colombia, la CEPAL los obtiene de IDEAM/UPME.

% ═════════════════════════════════════════════════════════════════════════════
\section{Simulador de Choques}
% ═════════════════════════════════════════════════════════════════════════════

Sobre las matrices publicadas se construyó un simulador que \emph{no} reconstruye
ni altera matrices, solo las usa como insumo. Considera dos familias de choque:
\begin{itemize}[leftmargin=1.6em,nosep]
  \item \textbf{Demanda} (cantidades de Leontief): $\Delta x = L\,\Delta y$,
        propaga efectos hacia atrás ante un cambio en la demanda final.
  \item \textbf{Oferta/costos} (precios de Ghosh):
        $\Delta x' = \Delta v'\,G$, propaga efectos hacia adelante ante un cambio
        en insumos primarios o valor agregado.
\end{itemize}
Las verificaciones algebraicas confirman que un choque unitario a la demanda del
sector $j$ reproduce exactamente la columna $j$ de $L$ y que el multiplicador de
demanda coincide con la suma de columna de $L$.

% ═════════════════════════════════════════════════════════════════════════════
\section{Reglas de Colaboración y Trabajo Futuro}
% ═════════════════════════════════════════════════════════════════════════════

\paragraph{Reglas para no romper la trazabilidad.}
\begin{itemize}[leftmargin=1.6em,nosep]
  \item No reemplazar una matriz publicada sin regenerar validación y auditoría.
  \item No borrar sectores por diagonal cero.
  \item No mezclar matrices directas con reconstruidas sin marcar el tipo.
  \item No presentar ajustes generales de $Z$ como cifras oficiales.
\end{itemize}

\paragraph{Trabajo futuro.}
\begin{itemize}[leftmargin=1.6em,nosep]
  \item Uruguay 2017: conseguir demanda final fuente completa para resolver la
        alerta de demanda final negativa.
  \item Extender Brasil y Uruguay a los años más recientes según se publiquen
        COU/TRU compatibles.
  \item Incorporar vectores de empleo y de emisiones por sector para activar los
        multiplicadores de empleo y la extensión ambiental.
\end{itemize}

\end{document}
